% Geodesique entre deux points et distance geodesique, comparees a une autre
% courbe tracee sur la variete.
% Utilise dans la \cref{sec:context-riemann-notion}.
% Styles et briques communs aux figures de geometrie riemannienne du
% \cref{ch:context-riemann}. Ce fichier peut etre \input plusieurs fois.
\definecolor{rmanifold}{RGB}{90,90,90}
\definecolor{rtangent}{RGB}{50,110,120}
\definecolor{rgeodesic}{RGB}{40,70,140}
\definecolor{rvector}{RGB}{175,95,45}
\tikzset{
  riemann/.style={
    >={Stealth[length=4.5pt,width=3.5pt]},
    line join=round,
    font=\footnotesize,
  },
  % La nappe qui represente la variete, centree sur l'origine.
  rsurface/.pic={
    \fill[rmanifold!7]
      (-3.2,0.15) .. controls (-2.95,1.75) and (-1.5,2.35) .. (0,2.35)
      .. controls (1.5,2.35) and (2.95,1.7) .. (3.2,0.15)
      .. controls (1.9,-0.7) and (-1.9,-0.7) .. (-3.2,0.15) -- cycle;
    \draw[rmanifold!70, thick]
      (-3.2,0.15) .. controls (-2.95,1.75) and (-1.5,2.35) .. (0,2.35)
      .. controls (1.5,2.35) and (2.95,1.7) .. (3.2,0.15);
    \draw[rmanifold!70, thick]
      (-3.2,0.15) .. controls (-1.9,-0.7) and (1.9,-0.7) .. (3.2,0.15);
  },
  % Le plan tangent, a placer en un point donne via un shift.
  rplane/.pic={
    \fill[rtangent!10, draw=rtangent!65]
      (-1.35,-0.42) -- (1.55,-0.88) -- (2.05,0.22) -- (-0.85,0.68) -- cycle;
  },
}

\begin{tikzpicture}[riemann]
  \coordinate (P) at (-1.95,1.05);
  \coordinate (Q) at ( 1.85,1.35);

  \path (0,0) pic {rsurface};
  \node[rmanifold!85] at (2.65,0.55) {$\Manifold$};

  % une courbe quelconque joignant les deux points
  \draw[rmanifold!60, thick, dashed]
    (P) .. controls (-1.20,-0.15) and (1.05,-0.05) .. (Q);
  \node[rmanifold, anchor=north, font=\scriptsize] at (-0.15,0.05)
    {$L(\geod')>\geodist(\param_0,\param_1)$};

  % la geodesique
  \draw[rgeodesic, very thick]
    (P) .. controls (-0.85,2.05) and (0.75,2.15) .. (Q);
  \node[rgeodesic, anchor=east, inner sep=3pt] at (-1.35,1.72) {$\geod$};
  \node[rgeodesic!85, anchor=south, font=\scriptsize] at (0.05,2.06)
    {$L(\geod)=\geodist(\param_0,\param_1)$};

  \fill (P) circle (1.7pt);
  \node[anchor=east, inner sep=3pt] at (P) {$\param_0$};
  \fill (Q) circle (1.7pt);
  \node[anchor=west, inner sep=3pt] at (Q) {$\param_1$};
\end{tikzpicture}
